\section{The Sharp Universal Gap}\label{app:universal}
We prove Theorem~\ref{thm:universal} and the program comparison
\begin{equation}
 \Phi_n(d,\A_{\rm fixed}(b))
 \le \Phi_{n/K_{m,b}}(d,\A_{\rm branch}(b)).
 \label{eq:programsimulation}
\end{equation}
Here the program is defined for any positive real sample-size argument.
All scores in the construction below depend only on the queried feature values; the estimator subsequently multiplies them by $Y$. Theorems~\ref{thm:finite} and~\ref{thm:asym} already compare these constructions with the larger class of history-dependent protocols, including branching choices based on the current label.

\subsection{Simulating One Tree without Branching}
Let $\tau$ be a deterministic tree of depth at most $b$ on $m\ge b$ coordinates, and let $H_\tau(x)$ be a score constant on each of its leaves. A repeated query on a path can be removed because its answer is already known. At every leaf with fewer than $b$ distinct queries, append unused coordinates in increasing index order until the path length is $b$. Copy the original score to all descendants. For a tree with no queries, start with any coordinate. The padded tree has exactly $b$ distinct queries on every path and the same score at every $x$.

Let $S_\tau(x)$ be its realized set of coordinates and let $\mathcal S_\tau$ be the collection of distinct such sets. Put $L_\tau=|\mathcal S_\tau|$. All sets contain the root, giving $L_\tau\le\binom{m-1}{b-1}$. The first $b-1$ answers fix all queried coordinates, even though the final answer can change the leaf score. Thus $L_\tau\le2^{b-1}$, proving $L_\tau\le K_{m,b}$.

Independently of $X,Y$, choose $S$ uniformly from $\mathcal S_\tau$ and observe $X_S$. Run the padded tree until a queried coordinate is unavailable or a leaf is reached. Let $M_{\tau,S}(X_S)$ be the indicator of reaching a leaf. This event is measurable from the fixed-form transcript. For each full profile $x$,
\[
 M_{\tau,S}(x_S)=1\quad\Longleftrightarrow\quad S=S_\tau(x).
\]
Indeed equality allows every successive query to be answered; conversely completion means that all $b$ distinct coordinates of the actual path lie in the chosen set of size $b$. Define the simulated score to be $\widetilde H_{\tau,S}=L_\tau H_\tau$ on completion and zero otherwise. Then, pointwise in $x$,
\begin{equation}
 \E_S[\widetilde H_{\tau,S}(x_S)]=H_\tau(x),\qquad
 \E_S[\widetilde H_{\tau,S}(x_S)^2]=L_\tau H_\tau(x)^2.
 \label{eq:simulate-moments}
\end{equation}

\subsection{Mixtures and the Finite-Sample Program}
For any feasible branching design $(q,g)$ take $H_\tau=g_\tau/q_\tau$ on its positive-mass trees. Draw the tree with probability $q_\tau$, then draw its path set as above. The resulting fixed-form procedure has conditional mean $s_g(x)$ and conditional second moment at most
\[
 K_{m,b}\sum_\tau g_\tau(x)^2/q_\tau.
\]
The same coordinate set can arise from several tree tags. To represent the construction using only the untagged fixed-mask menu, give $S$ mass $q'_S=\sum_{\tau:S\in\mathcal S_\tau}q_\tau/L_\tau$ and replace its score by the weighted average of the tagged scores, with weights $(q_\tau/L_\tau)/q'_S$. The tag is drawn before the person, so these weights do not depend on the population or the answers. Writing $p_{\tau,S}=q_\tau/L_\tau$ and $\overline H_S=\sum_\tau p_{\tau,S}\widetilde H_{\tau,S}/q'_S$, the reduction satisfies, pointwise,
\[
 \sum_Sq'_S\overline H_S=\sum_{\tau,S}p_{\tau,S}\widetilde H_{\tau,S},
 \qquad
 \sum_Sq'_S\overline H_S^2\le\sum_{\tau,S}p_{\tau,S}\widetilde H_{\tau,S}^2.
\]
This is conditional Jensen averaging over the tag given the fixed set. Zero-mass forms are discarded.

For an unbiased branching design this gives $V_{\rm fixed}\le K_{m,b}V_{\rm branch}$ after taking infima. Fixed forms are available as branching trees, so the reverse comparison $V_{\rm fixed}\ge V_{\rm branch}$ holds. Moreover, Jensen gives $V_{\rm branch}\ge\|d\|_\infty^2>0$. For a possibly biased design its residual $d-s_g$ is unchanged and its second moment is multiplied by at most $K_{m,b}$. Applying the same comparison to \eqref{eq:program} gives \eqref{eq:programsimulation}. No attainment assumption is needed: apply the construction to designs approaching the infimum.

\subsection{The Address Function Attains the Bound}
Let $\kappa=b-1$, $D=2^\kappa$ and $m=\kappa+D$. Identify the $\kappa$ address bits with an index $A\in\{1,\ldots,D\}$ in a fixed binary order, and let the remaining bits be $Z_1,\ldots,Z_D$. Set $d=Z_A$. This Boolean contrast is realized by the two rules $p_1=(1+d)/2$, $p_0=(1-d)/2$. Its polynomial expression is
\[
 d(A,Z)=\sum_{j=1}^D Z_j\mathbf1\{A=j\}.
\]
Each address indicator has degree $\kappa$, so $\deg(d)\le \kappa+1=b$ and both menus identify the contrast.

Under the uniform profile law $\U$, fix a mask with $s$ address bits and a set $R$ of $r$ observed data indices. For a realized partial address $\alpha$, let $C(\alpha)$ be its $2^{\kappa-s}$ consistent indices. Conditional independence and the centered unobserved data bits give
\[
 \E_\U[d\mid X_T]
      =2^{-(\kappa-s)}\sum_{j\in R\cap C(\alpha)} Z_j.
\]
The cross terms vanish on averaging over the observed independent signs. Each $j\in R$ is consistent with the random partial address with probability $2^{-s}$. Consequently
\[
 \|\E_\U[d\mid X_T]\|_2^2
   =2^{-2(\kappa-s)} r2^{-s}=r2^{s-2\kappa}.
\]
The value is zero if $r=0$. Otherwise $s\le \kappa+1-r$ and $r2^{1-r}\le1$ for every integer $r\ge1$, so it is at most $2^{-\kappa}=1/D$. Querying all $\kappa$ address bits and one data bit attains equality. Thus $J_\U(d)=1/D$ for fixed forms, whereas $\langle d,d\rangle_\U=1$. The dual formula \eqref{eq:vdual} gives $V_{\rm fixed}\ge D$. This is a minimax lower bound over all protocols in the model, not only over unbiased score estimators, by Theorem~\ref{thm:asym}.

For the upper bound choose $J$ uniformly from $\{1,\ldots,D\}$ before the interview, and observe all address bits and $Z_J$. The feature score $H=DZ_J\mathbf1\{J=A\}$ satisfies $\E_J[H\mid X]=d(X)$ and $\E_J[H^2\mid X]=D$ at every profile. Hence $V_{\rm fixed}=D$. Branching observes the address and then its data bit, giving feature score $d$ with second moment one. Jensen gives the matching lower bound, so $V_{\rm branch}=1$. In fact branching has the full-information risk $1/(1+\sqrt n)^2$ for every $n$ by the binary-experiment calculation in Appendix~\ref{app:sign}. The construction also covers $b=1$, with no address bits and one data bit. It proves the asserted equality for every $b$.

Finally, for either menu write $C=D$ or $C=1$ as appropriate. The uniform witness $h=d$ in \eqref{eq:dual} has denominator $1+n/C$ and numerator one. Shrinking the preceding unbiased score by $\gamma=n/(n+C)$ gives residual $(1-\gamma)d$ and constant second moment $\gamma^2 C$. The lower and upper bounds on the program therefore coincide:
\[
 \Phi_n=\frac{C}{n+C}.
\]
The finite-sample minimax risk is bounded using Theorem~\ref{thm:finite}; identifying its exact ratio for fixed and branching forms is not part of Theorem~\ref{thm:universal}.
